\documentclass[10pt,a4paper]{amsart}
\usepackage[margin=19mm]{geometry}
\usepackage{amsmath,amssymb,mathtools}
\usepackage[scr=rsfs,cal=euler]{mathalfa}
\usepackage{xfrac}
\usepackage[unicode,hidelinks,bookmarks=false]{hyperref}
\newcommand{\ie}{i.e.\ }
\newcommand{\cf}{cf.\ }
\newcommand{\resp}{respectively\ }
\newcommand{\Subsection}{\subsection}
\providecommand{\nobreakdash}{\nobreak\hbox{-}}
\setlength{\emergencystretch}{2em}
\multlinegap0pt
\usepackage[shortlabels]{enumitem}
\usepackage{amssymb}
\usepackage{xcolor}
\newtheorem{lemma}{Lemma}
\newtheorem{theorem}{Theorem}
\newcommand{\EE}{\mathcal{E}}
\newcommand{\NN}{\mathbb N}
\title{Fixed Points Theorems in Hausdorff M-distance Spaces}
\author{Vladyslav Babenko, Vira Babenko, Oleg Kovalenko}
\date{Source: arXiv:2111.13625v1 (CC BY 4.0)}
\begin{document}
\maketitle

\noindent\emph{Source and licence.} Fixed Points Theorems in Hausdorff M-distance Spaces. Version arXiv:2111.13625v1 (CC BY 4.0), distributed by arXiv under the Creative Commons Attribution 4.0 International licence, http://creativecommons.org/licenses/by/4.0/, as recorded in the arXiv licence metadata for this exact version and shown on the abstract page https://arxiv.org/abs/2111.13625v1. Full licence notice: \texttt{licenses/arxiv-2111.13625-CC-BY-4.0.txt}.\par\smallskip

\noindent\textit{Editorial scope.} Counted unit: the article's theorem th::M-KContraction (source0.tex lines 510--534). The article's complete original proof of it is delivered with it (source0.tex lines 536--595, one \texttt{proof} environment of 60 lines). No mathematics is written, repaired or re-derived: the statement and the proof are the article's own lines, in the article's own notation, and no retained line is substituted or re-flowed.  Retained uncounted as the source-local context the proof needs: lines 226--228; lines 372--374; lines 391--393.  Nothing was omitted from any retained span, so the package carries no omission record.  No retained line contains a citation, so the wrapper carries no bibliography and no bibliography entry.  No retained line contains a figure environment, an image-inclusion command or drawing code, so no image is delivered and the package declares no asset dependency.  Editorial formatting only, in addition: an AMS \texttt{amsart} preamble that restates the article's own declarations and macros used by the retained lines, namely the environments \texttt{lemma}, \texttt{theorem} on separate counters, together with the standard packages those lines need.  The retained environments print in this document as Lemma 1, Theorem 1 in order of appearance, whereas the article prints the same items with the numbering of its own full document.  The complete native source of the article as distributed in the arXiv e-print for this exact version is bundled unchanged as \texttt{sources/118918/source0.tex}.
\begin{equation}\label{zetaIsNullSq}
    \{\alpha_n\},\{\beta_n\}\in Z(M) \implies  \{\zeta (\alpha_n,\beta_n)\}\in Z(M).
\end{equation}

\begin{lemma}\label{l::econvergence}
If $Z(M) = Z_{\EE}(M)$, then for a sequence $\{ x_n\}\subset X$ and $x\in X$ one has that $x_n\to x$ {as} $n\to\infty$ if and only if for any $\varepsilon\in \EE$ there exists $N\in\NN$ such that for $n\geq N$ one has $d_{X,M}(x_n,x)< \varepsilon$.
\end{lemma}

 \begin{lemma}\label{l::eCauchySequence}
 A sequence $\{x_n\}$ in an $EDM$-distance space $(X, M, Z_{\EE}(M))$ is a Cauchy sequence if and only if for all $\varepsilon\in\EE$ there exists $N\in\NN$ such that $d_{X,M}(x_n,x_m)<\varepsilon$ for all $m\geq n\geq N$.
 \end{lemma}

\begin{theorem}\label{th::M-KContraction}
Let $(X,d_{X,M},Z_{\EE}(M))$ be an $\mathcal{H}$-complete $EDM_\zeta$-metric space with coordinate-wise non-decreasing $\zeta$ that satisfies property~\eqref{zetaIsNullSq} and 
$
\zeta(\varepsilon,\delta)\geq \varepsilon,\, \zeta(\varepsilon,\delta)\geq \delta \text{ for all } \varepsilon, \delta\in \EE.
$
Assume $f\colon X\to X$ satisfies the following property: for each $\varepsilon\in \EE$ there exists $\delta\in \EE$ such that for all $x,y\in X$ 
\begin{equation}\label{taylor2}          d_{X,M}(x,y)\leq \zeta( \delta, \varepsilon) \implies d_{X,M}(f(x),f(y))<\varepsilon.
\end{equation}
Then $f$ has a fixed point, provided one of the two sets of properties holds:
\begin{enumerate}[label = \Roman*]
\item\label{conditionsOnM}
\begin{enumerate}
    \item\label{taylorUniqueness1} If $x,y\in X$ and $\alpha, \beta\in M_+$ are such that $d_{X,M}(x,y)< \alpha+\beta$, then there exists $z\in X$ such that $d_{X,M}(x,z) < \alpha$ and $d_{X,M}(z,y) < \beta$;
    \item\label{taylorUniqueness2} For arbitrary $\varepsilon\in \EE$ and $x,y\in X$ there exists $n\in\NN$ such that $d_{X,M}(x,y)< n\cdot \varepsilon$;
\end{enumerate}
\item \label{conditionsOnF}
\begin{enumerate}
    \item  There exists $x\in X$ such that 
\begin{equation}\label{taylor1}
  \{  d_{X,M}(f^n(x),f^{n+1}(x))\}\in Z_{\EE}(M).
\end{equation} 
\end{enumerate}
\end{enumerate}
If conditions~\ref{conditionsOnM} hold, then the fixed point is unique.
\end{theorem}

\begin{proof}
First of all observe that $f$ is non-expansive in the following sense: for arbitrary $\varepsilon\in\EE$,
\begin{equation}\label{fIsNonExpanding}
    d_{X,M}(x,y)< \varepsilon\implies d_{X,M}(f(x),f(y))<\varepsilon.
\end{equation}
Indeed, let $\delta\in \EE$ be chosen according to condition~\eqref{taylor2}. Since  $d_{X,M}(x,y)<\varepsilon \leq \zeta(\delta,\varepsilon  )$, we obtain $d_{X,M}(f(x),f(y))<\varepsilon$ as required. 

Next we show that if conditions~\ref{conditionsOnM} hold, then for arbitrary $a,b\in X$, 
\begin{equation}\label{taylor0}
  \{  d_{X,M}(f^n(a),f^{n}(b))\}\in Z_{\EE}(M).
\end{equation}

Let $\varepsilon\in\EE$ be arbitrary, $\delta\in\EE$ be chosen according to~\eqref{taylor2} and $n$ be chosen according to property~\eqref{taylorUniqueness2}, so that $d_{X,M}(a,b)<n\cdot \delta$. Due to Lemma~\ref{l::econvergence}, inclusion~\eqref{taylor0} is implied by \begin{equation}\label{taylorUniquenessOfFixedPoint}
d_{X,M}(f^m(a), f^m(b))< \varepsilon\text{ for all } m\geq n.
\end{equation}
In order to prove~\eqref{taylorUniquenessOfFixedPoint}, it is sufficient to show that 
\begin{equation}\label{inductionStatement}
d_{X,M}(a,b)<n\cdot \delta\implies d_{X,M}(f^n(a), f^n(b))< \varepsilon
\end{equation}
in virtue of~\eqref{fIsNonExpanding}. We proceed by induction on $n\in\NN$. If $n=1$, then 
$
d_{X,M}(a,b) < \delta \leq \zeta(\delta, \varepsilon),
$
which implies~\eqref{inductionStatement} with $n = 1$, due to~\eqref{taylor2}.
Assume that it holds for some $n\in\NN$ and 
$$
d_{X,M}(a,b)< (n+1)\delta = \delta + n\delta.
$$
In view of property~\eqref{taylorUniqueness1}, there exists $c\in X$ such that $d_{X,M}(a,c)<  \delta$ and $d_{X,M}(c,b) < n\delta$. In virtue of the inductive assumption, $d_{X,M}(f^n(c), f^n(b))< \varepsilon$, and  inequality~\eqref{fIsNonExpanding} implies that  $d_{X,M} (f^n(a),f^n(c))< \delta$. Applying the triangle inequality, $
d_{X,M} (f^n(a),f^n(b))
\leq \zeta( \delta, \varepsilon)
$
and thus by~\eqref{taylor2}, $d_{X,M} (f^{n+1}(a),f^{n+1}(b))< \varepsilon$, as desired.

Applying~\eqref{taylor0} to $a = x$ and $b = f(x)$ with arbitrary $x\in X$, we obtain~\eqref{taylor1}. Hence in both cases when either conditions~\ref{conditionsOnM} or~\ref{conditionsOnF} hold, there is $x\in X$ such that~\eqref{taylor1} holds.

Next we show that  $\{f^n(x)\}$ is a Cauchy sequence.  For a fixed $\varepsilon\in \EE$ choose $\delta\in \EE$ according to condition~\eqref{taylor2}. Due to~\eqref{taylor1} and Lemma~\ref{l::econvergence}, there exists  $n_0$ such that for all $n\ge n_0$ one has
\begin{equation}\label{taylor1Conseq}
d_{X,M}(f^n(x),f^{n+1}(x))<\delta \text{ and } d_{X,M}(f^n(x),f^{n+1}(x))<\varepsilon.  
\end{equation}
Let $n> n_0$ be fixed. We show by induction on $k$ that for all $k\in \NN$  one has
\begin{equation}\label{fundamentalityCondition}
d_{X,M}(f^n(x),f^{n+k}(x))<\varepsilon.
\end{equation}
For $k=1$ inequality~\eqref{fundamentalityCondition} holds due to~\eqref{taylor1Conseq}. Assume that it is true for some $k \geq 1$. Then due to~\eqref{taylor1Conseq}, $d_{X,M}(f^{n-1}(x),f^{n}(x))<\delta$, and $d_{X,M}(f^n(x),f^{n+k}(x))<\varepsilon$ by the inductive assumption. Due to the triangle inequality, $d_{X,M}(f^{n-1}(x),f^{n+k}(x))\leq \zeta(\delta,\varepsilon)$, and hence in virtue of condition~\eqref{taylor2}, 
$
d_{X,M}(f^n(x),f^{n+k+1}(x))<\varepsilon,
$
which completes the induction step. Hence $\{f^n(x)\}$ is a Cauchy sequence due to Lemma~\ref{l::eCauchySequence}.

Completeness of the space implies existence of a point $\overline{x}\in X$ such that 
\begin{equation}\label{orbitLimit}
    f^n(x)\to\overline{x} \text{ as } n\to\infty.
\end{equation}
Property~\eqref{fIsNonExpanding} implies that $f$ is continuous, i.e. if $x_n\to x$, then $f(x_n)\to f(x)$. 
From~\eqref{orbitLimit} and continuity of $f$ we obtain $f(f^n(x))\to f(\overline{x})$. Uniqueness of the limit implies
$f(\overline{x})=\overline{x}$.

Finally, we prove the uniqueness of the fixed point  in the case, when conditions~\ref{conditionsOnM} hold. Let $a,b$ be two fixed points of $f$. Then for all $m\in\NN$,  $d_{X,M}(a,b) = d_{X,M}(f^m(a),f^m(b))$ and hence~\eqref{taylor0} implies that $\{d_{X,M}(a,b)\}\in Z(M)$, thus $a=b$, as desired.
\end{proof}

\end{document}
